\documentclass[11pt]{article}
\usepackage[a4paper,margin=1in]{geometry}
\usepackage{amsmath,amssymb,amsthm}
\usepackage{hyperref}

\newcommand{\dd}{\mathrm{d}}
\newcommand{\cD}{\mathcal{D}}
\newcommand{\cH}{\mathcal{H}}
\newcommand{\cQ}{\mathcal{Q}}
\newcommand{\avg}[1]{\langle #1\rangle}
\newcommand{\jap}[1]{\langle #1\rangle}

\newtheorem{theorem}{Theorem}
\newtheorem{lemma}[theorem]{Lemma}
\newtheorem{corollary}[theorem]{Corollary}
\theoremstyle{remark}
\newtheorem{remark}[theorem]{Remark}

\title{Solvability of the Chapman--Enskog problems and of the stationary
problem for the linearized Landau--Coulomb operator}
\author{Working note}
\date{16 November 2026}

\begin{document}
\sloppy
\maketitle

\noindent\emph{Summary.} The note \texttt{landau\_3d\_closure.tex} proved its
theorems for the Landau operator under two assumptions: (R), solvability of
two second-order Chapman--Enskog problems, and (S), existence of the
stationary kinetic solution together with its identification as the
quasi-static limit. This note proves (R), and the existence half of (S).
\begin{itemize}
  \item \textbf{Theorem~\ref{thm:solv}}: for the linearized like-particle
    Landau operator with Coulomb interaction, $\Lambda f=\sigma$ has a unique
    solution orthogonal to the conserved quantities whenever $\sigma$ is, and
    the solution inherits every polynomial moment of $\sigma$, with a fixed
    loss.
  \item \textbf{Corollary~\ref{cor:R}}: assumption (R) holds, and the
    constants $\kappa_{\mathrm F}$, $\eta$, $C_*$, $X_0$, $B$ are finite.
  \item \textbf{Theorem~\ref{thm:exist}}: for every $k\ne0$ the stationary
    problem has a solution, obtained variationally, for which every
    conclusion of the earlier theorems holds.
\end{itemize}
The proofs rest on published estimates for the Landau operator, listed in
Section~\ref{sec:cited} with what was and was not read. They do not reprove
those estimates. That the quasi-static limit of the time-dependent problem is given by this
stationary solution is proved in \texttt{landau\_quasistatic\_limit.tex}. No second
reader has checked the proofs.

\section{Setting and the cited estimates}
\label{sec:cited}

$M$ is the unit Maxwellian on $\mathbb{R}^3$, $\avg{\cdot}$ its average,
$w=\jap{v}=(1+|v|^2)^{1/2}$, $H$ is the span of $1,v_1,v_2,v_3,|v|^2$,
and $\Pi$ is the orthogonal projection onto $H$ in $L^2(M)$. $\cD$ is the
Dirichlet form of the linearized like-particle Landau operator with Coulomb
kernel,
\[
  \cD(g,h)=c_0\iint M(v)M(v')\,(\nabla g-\nabla'g')\cdot U(v-v')\,
  (\nabla h-\nabla'h')\,\dd v\,\dd v',\qquad
  U(z)=\frac{|z|^2I-zz^{\mathsf T}}{|z|^3}.
\]
The literature works with the Maxwellian $\mu=e^{-|v|^2}$, or with
$(2\pi)^{-3/2}e^{-|v|^2/2}$, and with $g=\sqrt\mu\,h\in L^2$ in place of
$h\in L^2(M)$; these differ from the present setting by a dilation of $v$ and
positive constants, which do not affect any statement below.

For $\theta\in\mathbb{R}$ let $\|h\|_{\sigma,\theta}$ denote Guo's weighted
norm $|g|_{\sigma,\theta}$ of $g=\sqrt M\,h$,
\[
  |g|^2_{\sigma,\theta}=\int w^{2\theta}\bigl[\sigma^{ij}\partial_ig\,\partial_jg
  +\sigma^{ij}v_iv_j\,g^2\bigr]\dd v,\qquad \sigma^{ij}=U^{ij}*\mu,
\]
$\cH_\theta=\{h:\|h\|_{\sigma,\theta}<\infty\}$, $\cH=\cH_0$,
$\|\cdot\|_\sigma=\|\cdot\|_{\sigma,0}$. The following are used.

\begin{enumerate}
  \item[(C1)] \emph{Structure of $\sigma$.} $\sigma^{ij}(v)$ has eigenvalue
    $\lambda_1(v)\sim(1+|v|)^{-3}$ in the direction of $v$ and
    $\lambda_2(v)\sim(1+|v|)^{-1}$ on $v^\perp$
    \cite[Lemma~2.3]{KGH}.
  \item[(C2)] \emph{Lower bound on the norm.} For every $\theta$,
    $\int M\,w^{2\theta-1}h^2\,\dd v\le C_\theta\|h\|^2_{\sigma,\theta}$,
    together with the corresponding gradient terms
    \cite[Lemma~2.5]{KGH}.
  \item[(C3)] \emph{Coercivity.} $\cD[h]\ge\delta\,\|h-\Pi h\|^2_\sigma$ for
    some $\delta>0$ \cite[Lemma~2.6]{KGH}. A constructive version, with the
    slightly weaker norm
    $\|h-\Pi h\|^2_{H^1(\jap v^{-3}M)}+\|\jap v^{-1/2}(h-\Pi h)\|^2_{L^2(M)}$,
    is \cite[Theorem~1.3]{Mouhot}.
  \item[(C4)] \emph{Weighted estimate.} For every $\theta\in\mathbb{R}$ there
    are a constant $C_\theta$ and a positive Gaussian $\rho$ such that
    \[
      \cD(h,w^{2\theta}h)\ \ge\ \tfrac12\|h\|^2_{\sigma,\theta}
      -C_\theta\int M\rho\,h^2\,\dd v .
    \]
    This follows from the two inequalities for the operators $A$ and $K$ in
    \cite[Lemma~2.7]{KGH}, with the remainder $C_\delta|\mu g|_2^2$; the
    same estimate with a compactly supported remainder is
    \cite[Lemma~9]{StrainGuo}.
  \item[(C5)] \emph{Continuity.}
    $|\cD(h_1,w^{2\theta}h_2)|\le C_\theta\|h_1\|_{\sigma,\theta}\|h_2\|_{\sigma,\theta}$.
\end{enumerate}
(C5) is not quoted from a single lemma. With $L=-A-K$ in the notation of
\cite{KGH}, the $K$ part and the $\partial_i\sigma^i$ part are bounded by
\cite[Lemmas~2.6, 2.7]{KGH}; the principal part of $A$ is bounded by the
Cauchy--Schwarz inequality for the non-negative matrix $\sigma$; and the
term in which the derivative falls on the weight is bounded by
$C\int w^{2\theta-1}\lambda_1|P_v\nabla g_1||g_2|$, which is controlled by
(C1) and (C2) because $w^{-2}\lambda_1\lesssim w^{-1}$.

\paragraph{What was read.} \cite{Mouhot}, \cite{StrainGuo} and the cited
lemmas of \cite{KGH} were read in full text. The lemmas of \cite{KGH} are
restatements of Lemma~3, Corollary~1, Lemma~5 and Lemma~6 of Guo
\cite{Guo02}, which was not available and was not consulted directly. That
the linearized Landau--Coulomb operator has no spectral gap in $L^2(M)$, so
that these weighted estimates are needed, is stated in \cite{CM}.

\paragraph{Consequences used below.} By (C3) and (C5), $\cD$ is a bounded
form on $\cH$, vanishes exactly on $H$, and is an inner product equivalent to
$\|\cdot\|^2_\sigma$ on $\cH^\perp=\{h\in\cH:\Pi h=0\}$. Polynomials lie in
every $\cH_\theta$. By (C2), for a polynomial $p$ and $h\in\cH$,
$|\avg{ph}|\le\avg{w\,p^2}^{1/2}\avg{w^{-1}h^2}^{1/2}<\infty$.

\section{Solvability with moments}

\begin{theorem}\label{thm:solv}
Let $\sigma$ be orthogonal to $H$ in $L^2(M)$ with
$\avg{w^{2\theta+1}\sigma^2}<\infty$ for some $\theta\ge0$. Then there is a
unique $f\in\cH^\perp$ with
\begin{equation}
  \cD(f,\varphi)=\avg{\sigma\varphi}\qquad\text{for all }\varphi\in\cH ,
  \label{eq:weak}
\end{equation}
and $f\in\cH_\theta$, with
\[
  \|f\|_{\sigma,\theta}\le C_\theta\bigl(\avg{w^{2\theta+1}\sigma^2}^{1/2}
  +\|f\|_\sigma\bigr).
\]
In particular $\avg{w^{2\theta-1}f^2}<\infty$.
\end{theorem}

\begin{proof}
\emph{Existence and uniqueness.} For $\varphi\in\cH$,
$\avg{\sigma\varphi}=\avg{\sigma(\varphi-\Pi\varphi)}$, so by (C2)
\[
  |\avg{\sigma\varphi}|\le\avg{w\sigma^2}^{1/2}
  \avg{w^{-1}(\varphi-\Pi\varphi)^2}^{1/2}\le C\,\|\varphi-\Pi\varphi\|_\sigma .
\]
The Riesz theorem in $(\cH^\perp,\cD)$ gives a unique $f\in\cH^\perp$ with
\eqref{eq:weak} for $\varphi\in\cH^\perp$; both sides vanish for
$\varphi\in H$.

\emph{Moments.} On $\cH_\theta$ define
\[
  b(g,\varphi)=\cD(g,w^{2\theta}\varphi)+C_\theta\int M\rho\,g\,\varphi\,\dd v,
  \qquad
  \ell(\varphi)=\avg{\sigma\,w^{2\theta}\varphi}
  +C_\theta\int M\rho\,f\,\varphi\,\dd v .
\]
By (C5) $b$ is bounded, and by (C4) $b(g,g)\ge\tfrac12\|g\|^2_{\sigma,\theta}$.
By (C2),
\[
  |\avg{\sigma w^{2\theta}\varphi}|\le\avg{w^{2\theta+1}\sigma^2}^{1/2}
  \avg{w^{2\theta-1}\varphi^2}^{1/2}\le C\|\varphi\|_{\sigma,\theta},
\]
so $\ell$ is bounded. The Lax--Milgram theorem gives $f_\theta\in\cH_\theta$ with
$b(f_\theta,\varphi)=\ell(\varphi)$ for all $\varphi\in\cH_\theta$.

Let $d=f_\theta-f\in\cH$. Since $\theta\ge0$, multiplication by $w^{-2\theta}$
maps $\cH$ into $\cH_\theta$ (the derivative of the weight gains a factor
$w^{-1}$ and is handled as in (C5)). Take $\varphi=w^{-2\theta}d$:
\[
  \cD(f_\theta,d)+C_\theta\int M\rho\,w^{-2\theta}f_\theta\,d
  =\avg{\sigma d}+C_\theta\int M\rho\,w^{-2\theta}f\,d .
\]
By \eqref{eq:weak}, $\cD(f,d)=\avg{\sigma d}$. Subtracting,
$\cD[d]+C_\theta\int M\rho\,w^{-2\theta}d^2=0$, so $d=0$. Hence
$f=f_\theta\in\cH_\theta$, and the bound follows from
$\tfrac12\|f\|^2_{\sigma,\theta}\le b(f,f)=\ell(f)$.
\end{proof}

\begin{remark}
Each application loses a fixed amount of weight: the source needs
$w^{2\theta+1}$ and the solution is controlled with $w^{2\theta-1}$. This is
the trace of the missing spectral gap. It is harmless for sources with all
polynomial moments. For polynomial sources a stronger, pointwise statement
is in the literature: $\Lambda^{-1}$ of a polynomial decays, relative to the
Maxwellian, slower than any inverse Gaussian power
\cite[Lemma~6.2]{DYY}.
\end{remark}

\section{Assumption (R)}

Write $\cH_\infty=\bigcap_{\theta\ge0}\cH_\theta$. By
Theorem~\ref{thm:solv}, if $\sigma\perp H$ has all polynomial moments then
$\Lambda^+\sigma:=f\in\cH_\infty\cap\cH^\perp$. Notation as in
\texttt{landau\_3d\_closure.tex}: $Q=\tfrac12v_x(|v|^2-5)$,
$\vartheta=v_\perp^2-2v_x^2$.

\begin{corollary}\label{cor:R}
For the linearized Landau--Coulomb operator:
\begin{enumerate}
  \item[(a)] $\Xi=\Lambda^+Q$ exists in $\cH_\infty$, is odd in $v_x$,
    satisfies $\avg{v_x\Xi}=0$, and
    $\kappa_{\mathrm F}=\cD[\Xi]\in(0,\infty)$; likewise
    $\Lambda^+\vartheta\in\cH_\infty$ and $\eta<\infty$.
  \item[(b)] $\chi_*=1-v_x^2+v_x\Xi/\kappa_{\mathrm F}$ is orthogonal to $H$
    and has all polynomial moments; $G=\Lambda^+\chi_*\in\cH_\infty$, and
    $C_*=\cD[G]$, $X_0=\avg{\vartheta G}$ are finite.
  \item[(c)] With $\tilde G=G+h$, $h\in\operatorname{span}\{1,|v|^2\}$ such
    that $\avg{\tilde G}=\avg{v_x^2\tilde G}=0$, the function $v_x\tilde G$ is
    orthogonal to $H$ and has all polynomial moments;
    $\Upsilon=\Lambda^+(v_x\tilde G)\in\cH_\infty$ and $B=\cD[\Upsilon]<\infty$.
\end{enumerate}
So assumption (R) holds, and (A3) holds with form domain $\cH$.
\end{corollary}

\begin{proof}
(a) $Q$ and $\vartheta$ are polynomials orthogonal to $H$. $\Lambda$
commutes with $v_x\to-v_x$, so by uniqueness $\Xi$ is odd;
$\Xi\in\cH^\perp$ gives $\avg{v_x\Xi}=0$; $\cD[\Xi]>0$ because $Q\ne0$.

(b) $\chi_*$ is even in $v_x$, hence orthogonal to $v_j$; orthogonality to
$1$ and $|v|^2$ was computed in the earlier note. Since $\Xi\in\cH_\infty$,
(C2) gives $\avg{w^{p}v_x^2\Xi^2}<\infty$ for every $p$.

(c) $v_x\tilde G$ is odd in $v_x$, so it is orthogonal to $1$, $|v|^2$,
$v_y$, $v_z$; and $\avg{v_x\cdot v_x\tilde G}=0$ by the choice of $h$.
Moments as in (b).
\end{proof}

\section{Existence of the stationary solution}

Fix $k\ne0$. Let $X^\infty$ be the space of $\varphi\in\cH_\infty$, even in
$v_x$ and axisymmetric, with $\avg{v_x^2\varphi}=0$. For
$\varphi\in X^\infty$, $v_x\varphi$ is orthogonal to $H$ with all moments, so
$f_\varphi=\Lambda^+(v_x\varphi)\in\cH_\infty$ and
\[
  \cQ(\varphi,\varphi')=\cD(\varphi,\varphi')+k^2\cD(f_\varphi,f_{\varphi'})
\]
is a symmetric non-negative form on $X^\infty$.

\begin{theorem}\label{thm:exist}
\begin{enumerate}
  \item[(i)] $\cQ$ is positive definite on $X^\infty$, and
    $\avg{\varphi}^2\le\bigl(1/(k^2\kappa_{\mathrm F})+C_*\bigr)\cQ[\varphi]$.
  \item[(ii)] The completion $\bar X$ of $(X^\infty,\cQ)$ is a space of
    functions contained in $\cH$, and there is a unique $g_e\in\bar X$ with
    $\cQ(g_e,\varphi)=\avg{\varphi}$ for all $\varphi\in\bar X$.
  \item[(iii)] There is $w\in\cH$, odd in $v_x$, such that the pair
    $(g_e,w)$ solves the stationary problem
    $ikv_x(g_e+iw)+\Lambda(g_e+iw)=1$ in the sense that
    \begin{align}
      \cD(g_e,\varphi)-k\avg{v_xw\,\varphi}&=\avg{\varphi}
        &&\text{for even }\varphi\in\cH_\infty, \label{eq:E1}\\
      \cD(w,\zeta)+k\avg{v_xg_e\,\zeta}&=0
        &&\text{for odd }\zeta\in\cH_\infty. \label{eq:E2}
    \end{align}
  \item[(iv)] For this solution, all conclusions of Theorem~1 (structure and
    two-sided bounds) and of the leading-coefficient theorem of
    \texttt{landau\_3d\_closure.tex} hold, with
    $n_1=\avg{g_e}=\cQ[g_e]$.
\end{enumerate}
\end{theorem}

\begin{proof}
(i) If $\cQ[\varphi]=0$ then $\varphi\in H$ and $f_\varphi\in H$, so
$v_x\varphi=\Lambda f_\varphi=0$ and $\varphi=0$. For the inequality let
$\psi_*=(v_x-\Xi/\kappa_{\mathrm F})/k\in\cH_\infty$, so
$1-kv_x\psi_*=\chi_*$ and $\cD[\psi_*]=1/(k^2\kappa_{\mathrm F})$. Then
$\avg{\varphi}=\avg{\varphi\chi_*}+k\avg{v_x\varphi\,\psi_*}
=\cD(G,\varphi)+k\,\cD(f_\varphi,\psi_*)$, and the Cauchy--Schwarz inequality
for $\cD$ and in $\mathbb{R}^2$ gives the claim.

(ii) Let $(\varphi_n)$ be $\cQ$-Cauchy. Then $\varphi_n-\Pi\varphi_n$
converges in $\cH^\perp$, $\avg{\varphi_n}$ converges by (i), and
$f_{\varphi_n}$ converges in $\cH^\perp$. Writing
$\Pi\varphi_n=a_n+b_n|v|^2$, the numbers $a_n+3b_n=\avg{\varphi_n}$ and
$a_n+5b_n=-\avg{v_x^2(\varphi_n-\Pi\varphi_n)}$ converge, the second because
$\avg{v_x^2\,\cdot\,}$ is continuous on $\cH$ by (C2). So $\varphi_n$
converges in $\cH$ to an even function with $\avg{v_x^2\varphi}=0$, and the
limit determines the Cauchy sequence up to $\cQ$-null sequences. The mean is
a bounded functional on $\bar X$ by (i); apply the Riesz theorem.

(iii) Let $\varphi_n\in X^\infty$ converge to $g_e$ in $\bar X$ and put
$w_0=-k\lim f_{\varphi_n}\in\cH^\perp$, $w=w_0+\beta v_x$. For odd
$\zeta\in\cH_\infty$, $\cD(f_{\varphi_n},\zeta)=\avg{v_x\varphi_n\zeta}$ and
$v_x\zeta$ has all moments, so in the limit
$\cD(w,\zeta)=-k\avg{v_xg_e\zeta}$, which is \eqref{eq:E2}. For
$\varphi\in X^\infty$,
$\cQ(g_e,\varphi)=\cD(g_e,\varphi)-k\,\cD(w_0,f_\varphi)
=\cD(g_e,\varphi)-k\avg{v_xw_0\,\varphi}$, and $\avg{v_x\cdot v_x\varphi}=0$,
so \eqref{eq:E1} holds on $X^\infty$ for every $\beta$. Every even
axisymmetric $\varphi\in\cH_\infty$ is $c\cdot1$ plus an element of
$X^\infty$, and \eqref{eq:E1} for $\varphi=1$ reads $-k\avg{v_xw}=1$, which
fixes $\beta$.

(iv) The proofs in \texttt{landau\_3d\_closure.tex} use only the following,
all of which hold here. From \eqref{eq:E1} with $\varphi=1$ and
$\varphi=\varphi_0=\tfrac12(5-|v|^2)$: $-k\avg{v_xw}=1$ and
$k\avg{Qw}=1$, hence $q_1=i/k$; and $\avg{v_x^2g_e}=0$ by (ii). The energy
identity is $\cQ[g_e]=\avg{g_e}$, with $k^2\lim\cD[f_{\varphi_n}]=\cD[w]$.
The second equation was tested only with $\psi_*$ and $\Xi$, which are in
$\cH_\infty$. The first was tested only with trial functions in $X^\infty$,
where it is the identity $\cQ(g_e,\varphi)=\avg{\varphi}$.
\end{proof}

\begin{remark}
Uniqueness is within $\bar X$. Uniqueness among all pairs in $\cH$ satisfying
\eqref{eq:E1}--\eqref{eq:E2} would require testing the equations with the
solution itself, and the pairing $\avg{v_xg_ew}$ is not controlled in the
energy space for the Coulomb operator. Weighted estimates for the stationary
kinetic solution would settle it; they are not proved here.
\end{remark}

\section{What remains open}

\begin{enumerate}
  \item \textbf{The quasi-static limit} is now proved in
    \texttt{landau\_quasistatic\_limit.tex}, which also shows that the
    stationary solution has all moments and is unique among weak solutions
    in $\cH_1$. Existence of the time-dependent weak solution in the
    weighted class is taken there from a standard method and is not written
    out.
  \item \textbf{Dependence on unread sources.} (C1)--(C4) are quoted from a
    paper that restates Guo's lemmas; the original was not consulted.
  \item \textbf{Scope.} Like-particle collisions only; the results on
    positivity at all wavenumbers and on the collisionless limit are
    unchanged and still open.
\end{enumerate}

\begin{thebibliography}{9}
\bibitem{Guo02} Y.~Guo, \emph{The Landau equation in a periodic box}, Comm.\
Math.\ Phys.\ \textbf{231} (2002) 391--434 (not consulted directly).
\bibitem{KGH} J.~Kim, Y.~Guo, H.~J. Hwang, \emph{A $L^2$ to $L^\infty$
approach for the Landau equation}, arXiv:1610.05346 (Lemmas 2.3, 2.5, 2.6,
2.7 read).
\bibitem{Mouhot} C.~Mouhot, \emph{Explicit coercivity estimates for the
linearized Boltzmann and Landau operators}, Comm.\ Partial Differential
Equations \textbf{31} (2006) 1321--1348, arXiv:math/0607538 (read).
\bibitem{StrainGuo} R.~M. Strain, Y.~Guo, \emph{Exponential decay for soft
potentials near Maxwellian}, Arch.\ Ration.\ Mech.\ Anal.\ \textbf{187}
(2008) 287--339 (read).
\bibitem{CM} K.~Carrapatoso, S.~Mischler, \emph{Landau equation for very
soft and Coulomb potentials near Maxwellians}, Ann.\ PDE \textbf{3} (2017),
arXiv:1512.01638 (Sections 1--3 read).
\bibitem{DYY} R.-J. Duan, D.~Yang, H.~Yu, \emph{Asymptotics toward viscous
contact waves for solutions of the Landau equation}, arXiv:2103.12929
(statement of Lemma 6.2 read; it cites Duan and Yu, Adv.\ Math.\ 362 (2020),
not consulted).
\end{thebibliography}

\end{document}
